\chapter{Conclusiones y Trabajo Futuro}
\label{cap:capitulo6}

\section{Conclusiones Generales}
\label{sec:conclusiones}
\textit{PoC exitoso demuestra viabilidad de gemelo digital web para 5G táctico con tecnologías modernas gratuitas.\\
Arquitectura tres capas adecuada para validación preliminar rápida.}

\section{Consecución de Objetivos}
\label{sec:objetivos_conseguidos}
\textit{OE1-OE5: todos conseguidos?? (modelado, interoperabilidad, backend, frontend, validación).}

\section{Competencias Adquiridas}
\label{sec:competencias}
\textit{Propagación radioeléctrica, desarrollo fullstack, GIS, gestión proyectos ágil, investigación técnica.}

\section{Aplicabilidad Práctica}
\label{sec:aplicabilidad}
\textit{Emergencias, eventos temporales, zonas remotas, formación universitaria.}

\section{Limitaciones del Trabajo}
\label{sec:limitaciones_trabajo}
\textit{Modelo simplificado, simulación offline, sin optimización ubicación, interfaz básica, sin telemetría real.}

\section{Líneas de Trabajo Futuro}
\label{sec:futuro}
\textit{Mejoras modelo: ray-tracing, obstáculos 3D, multi-frecuencia, indoor.\\
Arquitectura tiempo real: GPU, WebSocket streaming, cloud distribuido, caché Redis.\\
Funcionalidades: optimización automática TX (algoritmo genético), multi-antena, análisis interferencia, estimación capacidad.\\
Integración externa: equipos reales, telemetría vivo, GIS profesional (QGIS).\\
Usabilidad: PWA offline, app móvil nativa, colaboración multi-usuario.}

\section{Impacto y Difusión}
\label{sec:impacto}
\textit{Publicación IEEE, repositorio GitHub open-source, transferencia empresas, uso educativo URJC.}

\section{Reflexión Personal}
\label{sec:reflexion}

